\section{Simulations}
\label{sec:simulation}


We aim to compare TS-MR, Greedy-MR, and key baseline algorithms, via simulation.



\subsection{Synthetic datasets}
\label{sec:simulation_synthetic}



We conduct simulations on three representative synthetic examples.  We average simulation results over 100 independent runs for each of the examples. The results are shown in Figure~\ref{fig:synthetic}.




%The detailed experimental setups for these examples are presented in Appendix~\ref{sec:Synthetic_Environments}.
% \paragraph{Example 1. Stochastic linear bandit with uniformly and independently distributed actions.} This is a basic example of standard stochastic linear bandit problems without any extra structure. TS-Freq shows pessimistic regret due to the inflation of the posterior, while other algorithms in general perform well. 



% \paragraph{Example 2. Contextual bandits embedded in the linear bandit problem  \citep{abbasi2013online}.} In this setting, Greedy performs suboptimally due to a lack of exploration for some arms. Nevertheless, Greedy-MR outperforms both Greedy and OFUL by adaptively choosing OFUL actions only when it detects large regret proxy $\hat{\mu}_t$.  



% \paragraph{Example 3. Prior mean mismatch \citep{hamidi2020frequentist}.} This is an example in which LinTS is proven to incur linear Bayesian regret.  We see both LinTS and Greedy incur linear regrets as expected, while TS-MR and Greedy-MR, switch to OFUL adaptively to tackle this hard problem and achieve sublinear regret.


%\section{Details of Synthetic Environments}
%\label{sec:Synthetic_Environments}
\paragraph{Example 1. Stochastic linear bandit with uniformly and independently distributed actions.} We fix $d = 50$, and sample $\param\sim\text{Unif}(\{\theta \in \mathbb{R}^d|\|\theta\|= 10\}) $ on a sphere with fixed norm. At each time $t$, we generate 100 i.i.d. random actions sampled from $\text{Unif}(\cS_{d-1})$ to form $\cX_t$. This is a basic example of standard stochastic linear bandit problems without any extra structure. We set the threshold  $\mu = 8$ for TS-MR and Greedy-MR. TS-Freq shows pessimistic regret due to the inflation of the posterior, while other algorithms in general perform well. 



\paragraph{Example 2. Contextual bandits embedded in the linear bandit problem  \citep{abbasi2013online}.} We fix $d = 50$, and sample $\param\sim\text{Unif}(\{\theta \in \mathbb{R}^d|\|\theta\|= 70\}) $. At each time $t$, we first  generate a random vector $u_t \sim \cN(0, I_5)$ and let $x_{t,i}\in \RR^{50}$ be the vector whose $i$-th block of size 5 is a  copy of $u_t$, and other components are 0. Then $X_t = \{x_{t,i}\}_{i\in [10] }$ is an action set of size 10, sharing the same feature $u_t$ in different blocks. This problem is equivalent to a 10-armed contextual bandit. We set $\mu = 12$ for TS-MR and Greedy-MR.  In this setting, Greedy performs suboptimally due to a lack of exploration for some arms. Nevertheless, Greedy-MR outperforms both Greedy and OFUL by adaptively choosing OFUL actions only when it detects large regret proxy $\hat{\mu}_t$.  



\paragraph{Example 3. Prior mean mismatch \citep{hamidi2020frequentist}.} This is an example in which LinTS is shown to incur linear Bayesian regret. We sample $\theta^\star \sim \cN(m \mathbf{1}_{3d}, I_{3d})$ and fix the action set $\cX_t = \{0, x_a, x_b\}$ for all $t\in[T]$, where $x_a = - \sum_{i=1}^{d} e_i/\sqrt{3d},  \quad x_b = \sum_{i=11}^{3d} e_i\sqrt{3d} - \sum_{i=1}^{d}e_i /\sqrt{3d}$.
It is shown in \cite{hamidi2020frequentist} that, when LinTS takes a wrong prior mean as input,  it has a large probability to choose  $\tilde{x}_2 = 0$, conditioned on $\tilde{x}_1 = x_a$. Note that choosing the zero action brings no information update to LinTS, it suffers a linear Bayesian regret when trying to escape from the zero action. We let $m = 10$ and set $d = 10$, so the problem is a $30$-dimensional linear bandit. We set $\mu = 12$ for TS-MR and Greedy-MR. We see both LinTS and Greedy incur linear regrets as expected, while TS-MR and Greedy-MR, switch to OFUL adaptively to tackle this hard problem and achieve sublinear regret.
\begin{figure}[!htp]
	\centering
	\begin{subfigure}[b]{0.32\textwidth}
		\includegraphics[width=\textwidth]{figures/cumuregrets_ex1.jpg}
				\caption{example 1}
				\label{fig:ex1}
\end{subfigure}
	\begin{subfigure}[b]{0.32\textwidth}
				\includegraphics[width=\textwidth]{figures/cumuregrets_ex2.jpg}
						\caption{example 2}
						\label{fig:ex2}
\end{subfigure}
	\begin{subfigure}[b]{0.32\textwidth}
		\includegraphics[width=\textwidth]{figures/cumuregrets_ex3.jpg}
				\caption{example 3}
				\label{fig:ex3}
\end{subfigure}\\
	\begin{subfigure}[b]{0.32\textwidth}
		\includegraphics[width=\textwidth]{figures/oful_fraction_ex1.jpg} 
	\caption{example 1}
		\label{fig:oful_fraction_ex1}
\end{subfigure}
	\begin{subfigure}[b]{0.32\textwidth}
	\includegraphics[width=\textwidth]{figures/oful_fraction_ex2.jpg} 
	\caption{example 2}
		\label{fig:oful_fraction_ex2}
	\end{subfigure}
 	\begin{subfigure}[b]{0.32\textwidth}
	\includegraphics[width=\textwidth]{figures/oful_fraction_ex3.jpg} 
	\caption{example 3}
		\label{fig:oful_fraction_ex3}
	\end{subfigure}
	\caption{Simulation results on synthetic data. (a) - (c): Cumulative regret of TS-MR and Greedy-MR versus baseline algorithms. Shaded regions show $\pm 2$ SE of mean regret. (d) - (f): Fraction of OFUL actions in TS-MR and Greedy-MR.}
	\label{fig:synthetic}
\end{figure}

\subsection{Real-world datasets}
\label{sec:simulation_real_world}
We explore the performance of standard \POFUL{} algorithms and the proposed TS-MR and Greedy-MR algorithms on real-world datasets. We use three classification datasets from \href{https://www.openml.org/}{\textsc{OpenML}}: \href{https://www.openml.org/d/1466}{Cardiotocography}, \href{https://www.openml.org/d/375}{JapaneseVowels}, and \href{https://www.openml.org/d/36}{Segment}, representing healthcare, pattern recognition, and computer vision domains. Following \citet{bietti2021contextual, bastani2021mostly}, we convert these classification tasks to contextual bandit problems and embed them into linear bandit problems as in Example 2, Section \ref{sec:simulation_synthetic}. Each class becomes an action where the decision-maker receives a binary reward (1 for correct classification, 0 otherwise) plus Gaussian noise.

We plot the cumulative regret (averaged over 100 runs) for all algorithms.   Figure~\ref{fig:emp} shows that for all real-world datasets: OFUL and TS-Freq perform poorly due to their conservative exploration; LinTS and Greedy are achieving empirical success even though they don't have theoretical guarantees; TS-MR and Greedy-MR retain the desirable empirical performance of LinTS and Greedy, while enjoying the minimax optimal frequentist regret bound.

 \begin{figure}[!htp]
	\centering
	\begin{subfigure}[b]{0.32\textwidth}
		\includegraphics[width=\textwidth]{figures/cumuregrets_emp1.jpg} 
	\caption{Cardiotocography dataset}
		\label{fig:Cardiotocography}
\end{subfigure}
	\begin{subfigure}[b]{0.32\textwidth}
	\includegraphics[width=\textwidth]{figures/cumuregrets_emp2.jpg} 
	\caption{JapaneseVowels dataset}
		\label{fig:JapaneseVowels}
	\end{subfigure}
 	\begin{subfigure}[b]{0.32\textwidth}
	\includegraphics[width=\textwidth]{figures/cumuregrets_emp3.jpg} 
	\caption{Segment dataset}
		\label{fig:segment}
	\end{subfigure}\\
    \begin{subfigure}[b]{0.32\textwidth}
		\includegraphics[width=\textwidth]{figures/oful_fraction_emp1.jpg} 
	\caption{Cardiotocography dataset}
		\label{fig:Cardiotocography_frac}
\end{subfigure}
	\begin{subfigure}[b]{0.32\textwidth}
	\includegraphics[width=\textwidth]{figures/oful_fraction_emp2.jpg} 
	\caption{JapaneseVowels dataset}
		\label{fig:JapaneseVowels_frac}
	\end{subfigure}
 	\begin{subfigure}[b]{0.32\textwidth}
	\includegraphics[width=\textwidth]{figures/oful_fraction_emp3.jpg} 
	\caption{Segment dataset}
		\label{fig:segment_frac}
	\end{subfigure}
	\caption{Simulation results on real-world datasets. (a) - (c): Cumulative regret of all algorithms. Shaded regions show the $\pm 2$ SE of the mean regret. (d) - (f): Fraction of OFUL actions of TS-MR and Greedy-MR.}
	\label{fig:emp}
\end{figure}






\begin{remark}
Simulation results in Figures~\ref{fig:synthetic} and \ref{fig:emp} show OFUL actions are primarily used in the early stages. This indicates: (1) Greedy-MR and TS-MR implement OFUL actions only when necessary, maintaining a low OFUL fraction throughout most of the time horizon, substantially reducing computational cost; and (2) limited course-corrected exploration at the beginning efficiently remedies TS and Greedy in problematic instances.
\end{remark}
